% Section 04 - Génération d'attestations de réussite
\section{Génération d'attestations de réussite}

\subsection{Vue d'ensemble}

Le menu \menu{Générer les attestations} permet de créer des attestations officielles pour les étudiants ayant \textbf{validé} leur semestre (moyenne $\geq$ 10/20).

\attention{IMPORTANT : Seuls les étudiants avec une moyenne générale supérieure ou égale à 10/20 recevront une attestation. Les autres sont automatiquement exclus.}

\textbf{Caractéristiques principales :}
\begin{itemize}[leftmargin=*]
    \item Validation automatique des moyennes
    \item Thèmes entièrement personnalisables
    \item 6 préréglages professionnels
    \item Support bilingue français/anglais
    \item QR code chiffré compact
    \item Statistiques en temps réel
\end{itemize}

\subsection{Différences avec les relevés de notes}

\begin{table}[h]
\centering
\begin{tabular}{|l|p{5.5cm}|p{5.5cm}|}
\hline
\textbf{Critère} & \textbf{Relevés de notes} & \textbf{Attestations de réussite} \\
\hline
Public cible & Tous les étudiants & Étudiants admis uniquement ($\geq$ 10/20) \\
\hline
Contenu & Notes détaillées par EC, moyennes par UE & Résultat global (admis), pas de détail des notes \\
\hline
Personnalisation & Thème lié à la configuration & Thème indépendant et très personnalisable \\
\hline
Langues & Français uniquement & Bilingue FR/EN \\
\hline
QR Code & Chiffré AES-128-ECB & Chiffré compact \\
\hline
\end{tabular}
\caption{Comparaison relevés vs attestations}
\end{table}

\subsection{Préparation du fichier Excel}

\subsubsection{Colonnes obligatoires}

\begin{tcolorbox}[colback=blue!5!white,colframe=blue!75!black,title=Colonnes obligatoires]
\begin{longtable}{|l|p{9cm}|}
\hline
\textbf{Colonne} & \textbf{Description} \\
\hline
\endfirsthead
\hline
\textbf{Colonne} & \textbf{Description} \\
\hline
\endhead
\hline
\endfoot

\texttt{MATRICULE} & Numéro matricule unique \\
\hline
\texttt{NOM} & Nom de famille \\
\hline
\texttt{PRENOM} & Prénom(s) \\
\hline
\texttt{DATE\_NAISSANCE} & Date de naissance (format JJ/MM/AAAA ou DD/MM/YYYY) \\
\hline
\texttt{LIEU\_NAISSANCE} & Lieu de naissance \\
\hline
\texttt{ANNEE\_ACADEMIQUE} & Année académique (ex: 2023-2024) \\
\hline
\texttt{SEMESTRE} & Code du semestre (S1, S2, etc.) \\
\hline
\texttt{NIVEAU} & Niveau d'études (L1, L2, L3, M1, M2, etc.) \\
\hline
\texttt{DOMAINE} & Domaine d'études (ex: Sciences de la Santé) \\
\hline
\texttt{PARCOURS} & Parcours ou spécialité (ex: Médecine Générale) \\
\hline
\texttt{MOYENNE\_GENERALE} & Moyenne générale du semestre (sur 20) \\
\hline
\end{longtable}
\end{tcolorbox}

\subsubsection{Colonnes optionnelles pour support bilingue}

Si vous souhaitez générer des attestations en anglais, ajoutez les colonnes de traduction :

\begin{longtable}{|l|l|}
\hline
\textbf{Colonne française} & \textbf{Colonne anglaise} \\
\hline
\endfirsthead
\hline
\textbf{Colonne française} & \textbf{Colonne anglaise} \\
\hline
\endhead
\hline
\endfoot

\texttt{DOMAINE} & \texttt{DOMAINE\_EN} \\
\hline
\texttt{PARCOURS} & \texttt{PARCOURS\_EN} \\
\hline
\texttt{SPECIALITE} & \texttt{SPECIALITE\_EN} \\
\hline
\texttt{OPTION} & \texttt{OPTION\_EN} \\
\hline
\texttt{FINALITE} & \texttt{FINALITE\_EN} \\
\hline
\texttt{MENTION} & \texttt{MENTION\_EN} \\
\hline
\end{longtable}

\begin{figure}[h]
    \centering
    % \includegraphics[width=\textwidth]{images/excel_attestations.png}
    \fbox{\parbox{0.95\textwidth}{\centering [CAPTURE D'ÉCRAN : EXEMPLE FICHIER EXCEL POUR ATTESTATIONS]\\ \vspace{0.3cm}
    Montrant toutes les colonnes obligatoires + colonnes bilingues}}
    \caption{Structure du fichier Excel pour attestations de réussite}
\end{figure}

\subsection{Workflow complet de génération}

\subsubsection{Étape 1 : Accéder au menu}

\begin{enumerate}[leftmargin=*]
    \item Cliquez sur \menu{Générer les attestations} dans la barre latérale
    \item L'interface s'affiche avec 6 onglets :
    \begin{enumerate}
        \item \textbf{Générateur} : Import Excel et génération
        \item \textbf{Sélection} : Filtrage des étudiants éligibles
        \item \textbf{Préréglages} : Thèmes prédéfinis
        \item \textbf{Thème} : Personnalisation basique
        \item \textbf{Style Avancé} : Personnalisation typographique
        \item \textbf{Paramètres} : Options d'export et QR code
    \end{enumerate}
\end{enumerate}

\subsubsection{Étape 2 : Charger le fichier Excel}

\begin{enumerate}[leftmargin=*]
    \item Dans l'onglet \textbf{Générateur}, cliquez sur \bouton{Charger fichier Excel}
    \item Sélectionnez votre fichier contenant les données des étudiants
    \item L'application valide automatiquement :
    \begin{itemize}
        \item Présence des colonnes obligatoires
        \item Format de la colonne \texttt{MOYENNE\_GENERALE}
        \item Format des dates
    \end{itemize}
    \item Un message de confirmation s'affiche
\end{enumerate}

\subsubsection{Étape 3 : Analyse automatique de l'éligibilité}

Dès le chargement, l'application analyse automatiquement les moyennes :

\begin{tcolorbox}[colback=green!5!white,colframe=green!75!black,title=Statistiques affichées]
\textbf{Exemple d'affichage :}

\begin{itemize}[leftmargin=*]
    \item \textcolor{green}{\textbf{✓ 127 étudiants ADMIS}} (moyenne $\geq$ 10/20)
    \item \textcolor{red}{\textbf{✗ 23 étudiants AJOURNÉS}} (moyenne < 10/20)
    \item \textcolor{blue}{\textbf{Total : 150 étudiants}}
\end{itemize}

\textbf{Taux de réussite : 84.67\%}
\end{tcolorbox}

\begin{figure}[h]
    \centering
    % \includegraphics[width=0.7\textwidth]{images/stats_eligibilite.png}
    \fbox{\parbox{0.65\textwidth}{\centering [CAPTURE D'ÉCRAN : STATISTIQUES D'ÉLIGIBILITÉ]\\ \vspace{0.3cm}
    Graphique circulaire ou barres montrant admis/ajournés\\
    Pourcentages et nombres absolus}}
    \caption{Statistiques en temps réel des étudiants éligibles}
\end{figure}

\attention{Les étudiants ajournés n'apparaîtront pas dans la génération finale, même s'ils sont présents dans l'Excel.}

\subsubsection{Étape 4 : Choisir ou personnaliser un thème}

Vous avez deux options pour le design de vos attestations :

\paragraph{Option A : Utiliser un préréglage (rapide)}

\begin{enumerate}[leftmargin=*]
    \item Accédez à l'onglet \textbf{Préréglages}
    \item L'application propose 6 thèmes professionnels :
    \begin{itemize}
        \item \textbf{Classique Élégant} : Design sobre avec bleu marine et or
        \item \textbf{Moderne Minimaliste} : Design épuré avec gris et vert
        \item \textbf{Professionnel Académique} : Style traditionnel universitaire
        \item \textbf{Dynamique Coloré} : Couleurs vives et modernes
        \item \textbf{Excellence Prestige} : Design luxueux avec bordeaux et or
        \item \textbf{Nature Zen} : Tons verts et beiges apaisants
    \end{itemize}
    \item Cliquez sur \bouton{Aperçu} pour voir un exemple
    \item Cliquez sur \bouton{Appliquer ce thème} pour l'utiliser
\end{enumerate}

\begin{figure}[h]
    \centering
    % \includegraphics[width=\textwidth]{images/prereglages_themes.png}
    \fbox{\parbox{0.9\textwidth}{\centering [CAPTURE D'ÉCRAN : GALERIE DES PRÉRÉGLAGES]\\ \vspace{0.3cm}
    6 cartes affichant les aperçus des thèmes\\
    Boutons "Aperçu" et "Appliquer" sous chaque thème}}
    \caption{Galerie des 6 préréglages de thèmes disponibles}
\end{figure}

\paragraph{Option B : Personnaliser entièrement (avancé)}

\begin{enumerate}[leftmargin=*]
    \item Accédez à l'onglet \textbf{Thème}
    \item Configurez les éléments suivants :

    \subparagraph{Couleurs}
    \begin{itemize}
        \item \champ{Couleur primaire} : Couleur principale du design
        \item \champ{Couleur secondaire} : Couleur d'accentuation
        \item \champ{Couleur d'arrière-plan} : Fond de l'attestation
        \item \champ{Couleur du texte} : Texte principal
    \end{itemize}

    \subparagraph{En-tête}
    \begin{itemize}
        \item \champ{Afficher le logo} : Oui/Non
        \item \champ{Position du logo} : Gauche / Centre / Droite
        \item \champ{Taille du logo} : Petit / Moyen / Grand
        \item \champ{Afficher le nom de l'établissement} : Oui/Non
    \end{itemize}

    \subparagraph{Corps du document}
    \begin{itemize}
        \item \champ{Alignement du texte} : Gauche / Centre / Justifié
        \item \champ{Espacement entre lignes} : 1.0 / 1.5 / 2.0
        \item \champ{Afficher la mention} : Oui/Non
        \item \champ{Afficher la moyenne} : Oui/Non
    \end{itemize}

    \subparagraph{Pied de page}
    \begin{itemize}
        \item \champ{Signature} : Charger une image de signature
        \item \champ{Nom du signataire}
        \item \champ{Fonction du signataire}
        \item \champ{Date de délivrance} : Date du jour / Personnalisée
    \end{itemize}

    \item Prévisualisez en temps réel dans le panneau de droite
    \item Enregistrez votre thème personnalisé
\end{enumerate}

\begin{figure}[h]
    \centering
    % \includegraphics[width=\textwidth]{images/personnalisation_theme.png}
    \fbox{\parbox{0.9\textwidth}{\centering [CAPTURE D'ÉCRAN : INTERFACE DE PERSONNALISATION]\\ \vspace{0.3cm}
    Panneau gauche : Formulaire avec tous les réglages\\
    Panneau droit : Aperçu en temps réel de l'attestation}}
    \caption{Interface de personnalisation complète du thème}
\end{figure}

\subsubsection{Étape 5 : Configuration typographique avancée (optionnel)}

Pour un contrôle total sur la typographie :

\begin{enumerate}[leftmargin=*]
    \item Accédez à l'onglet \textbf{Style Avancé}
    \item Configurez :

    \paragraph{Polices de caractères}
    \begin{itemize}
        \item \champ{Famille de police pour titres} : Serif / Sans-serif / Monospace / Police personnalisée
        \item \champ{Famille de police pour le corps} : Idem
        \item \champ{Taille du titre principal} : 24pt à 48pt
        \item \champ{Taille des sous-titres} : 14pt à 24pt
        \item \champ{Taille du corps de texte} : 10pt à 16pt
    \end{itemize}

    \paragraph{Effets de texte}
    \begin{itemize}
        \item \champ{Style du titre} : Normal / Gras / Italique / Gras Italique
        \item \champ{Lettrage} : Normal / Petites capitales / Majuscules
        \item \champ{Ombrage du texte} : Oui/Non + couleur et décalage
    \end{itemize}

    \paragraph{Bordures et cadres}
    \begin{itemize}
        \item \champ{Bordure du document} : Oui/Non
        \item \champ{Style de bordure} : Simple / Double / Décorative
        \item \champ{Épaisseur} : 1px à 10px
        \item \champ{Couleur de bordure}
        \item \champ{Arrondi des coins} : 0px (carré) à 20px (très arrondi)
    \end{itemize}

    \item Les modifications sont visibles en temps réel dans l'aperçu
\end{enumerate}

\subsubsection{Étape 6 : Configurer les paramètres d'export}

\begin{enumerate}[leftmargin=*]
    \item Accédez à l'onglet \textbf{Paramètres}
    \item Configurez les options suivantes :

    \paragraph{Format d'exportation}
    \begin{description}[leftmargin=3.5cm]
        \item[ZIP (défaut)] Archive contenant tous les PDFs
        \item[Fichiers individuels] Téléchargement séparé de chaque attestation
        \item[PDF unique] Toutes les attestations fusionnées en un seul fichier
    \end{description}

    \paragraph{Compression}
    \begin{itemize}
        \item Activez \champ{Utiliser la compression} pour réduire la taille du ZIP
    \end{itemize}

    \paragraph{QR Code chiffré compact}
    \begin{itemize}
        \item Activez \champ{Ajouter un QR code chiffré}
        \item Le QR code contient :
        \begin{itemize}
            \item Matricule de l'étudiant
            \item Nom et prénom
            \item Semestre et année académique
            \item Hash de vérification
        \end{itemize}
        \item Chiffrement : AES-128-ECB
        \item Position : Coin inférieur droit par défaut
    \end{itemize}

    \paragraph{Langue}
    \begin{itemize}
        \item \champ{Langue de génération} : Français / Anglais
        \item Si Anglais sélectionné, les colonnes \texttt{*\_EN} de l'Excel seront utilisées
    \end{itemize}
\end{enumerate}

\begin{figure}[h]
    \centering
    % \includegraphics[width=0.8\textwidth]{images/parametres_attestations.png}
    \fbox{\parbox{0.75\textwidth}{\centering [CAPTURE D'ÉCRAN : ONGLET PARAMÈTRES]\\ \vspace{0.3cm}
    Format : ● ZIP  ○ Fichiers individuels  ○ PDF unique\\
    ☑ Utiliser la compression\\
    ☑ Ajouter un QR code chiffré\\
    Langue : [Dropdown: Français]}}
    \caption{Configuration des paramètres d'exportation}
\end{figure}

\subsubsection{Étape 7 : Filtrer les étudiants (optionnel)}

Par défaut, des attestations seront générées pour TOUS les étudiants admis. Pour restreindre :

\begin{enumerate}[leftmargin=*]
    \item Accédez à l'onglet \textbf{Sélection}
    \item Vous voyez la liste complète des étudiants éligibles (moyenne $\geq$ 10)
    \item Trois options de filtrage :

    \paragraph{Filtrer par niveau}
    \begin{itemize}
        \item Cochez les niveaux souhaités : L1, L2, L3, M1, M2, etc.
    \end{itemize}

    \paragraph{Filtrer par semestre}
    \begin{itemize}
        \item Cochez les semestres souhaités : S1, S2, S3, S4, S5, S6
    \end{itemize}

    \paragraph{Sélection manuelle}
    \begin{itemize}
        \item Cochez individuellement les étudiants
        \item Utilisez la barre de recherche pour trouver rapidement un étudiant
    \end{itemize}

    \item Un compteur affiche : \texttt{45/127 étudiants sélectionnés}
\end{enumerate}

\begin{figure}[h]
    \centering
    % \includegraphics[width=\textwidth]{images/selection_attestations.png}
    \fbox{\parbox{0.9\textwidth}{\centering [CAPTURE D'ÉCRAN : ONGLET SÉLECTION]\\ \vspace{0.3cm}
    Liste des étudiants admis avec cases à cocher\\
    Filtres par niveau et semestre en haut\\
    Barre de recherche\\
    Compteur "45/127 sélectionnés"}}
    \caption{Filtrage et sélection des étudiants pour attestations}
\end{figure}

\subsubsection{Étape 8 : Générer les attestations}

\begin{enumerate}[leftmargin=*]
    \item Retournez sur l'onglet \textbf{Générateur}
    \item Vérifiez le récapitulatif affiché :
    \begin{itemize}
        \item Nombre d'attestations à générer
        \item Thème sélectionné
        \item Format d'exportation
        \item Options activées (compression, QR code, etc.)
    \end{itemize}
    \item Cliquez sur \bouton{Générer les attestations}
    \item Une barre de progression s'affiche
    \item Attendez la fin du traitement
\end{enumerate}

\begin{figure}[h]
    \centering
    % \includegraphics[width=0.8\textwidth]{images/generation_attestations.png}
    \fbox{\parbox{0.75\textwidth}{\centering [CAPTURE D'ÉCRAN : GÉNÉRATION EN COURS]\\ \vspace{0.3cm}
    Génération de 127 attestations...\\
    ▓▓▓▓▓▓▓▓▓░░░░░░  63/127 (49\%)\\
    Temps restant estimé : 1 min 30 s}}
    \caption{Progression de la génération des attestations}
\end{figure}

\subsubsection{Étape 9 : Téléchargement et vérification}

Une fois la génération terminée :

\begin{itemize}[leftmargin=*]
    \item Un message de succès s'affiche : \texttt{✓ 127 attestations générées avec succès}
    \item Le fichier ou les fichiers sont téléchargés automatiquement
    \item Ouvrez quelques attestations pour vérifier :
    \begin{itemize}
        \item Les informations de l'étudiant sont correctes
        \item Le thème est appliqué correctement
        \item Le QR code est présent (si activé)
        \item La mise en page est satisfaisante
    \end{itemize}
\end{itemize}

\subsection{Export et import de thèmes personnalisés}

Pour réutiliser vos thèmes ou les partager :

\subsubsection{Exporter un thème}

\begin{enumerate}[leftmargin=*]
    \item Dans l'onglet \textbf{Thème}, configurez votre thème personnalisé
    \item Cliquez sur \bouton{Exporter le thème}
    \item Donnez un nom au thème : "Mon Thème Personnalisé"
    \item Un fichier JSON est téléchargé : \texttt{theme\_mon\_theme\_personnalise.json}
    \item Conservez ce fichier en lieu sûr
\end{enumerate}

\subsubsection{Importer un thème}

\begin{enumerate}[leftmargin=*]
    \item Dans l'onglet \textbf{Thème}, cliquez sur \bouton{Importer un thème}
    \item Sélectionnez le fichier JSON du thème
    \item Le thème est chargé et appliqué immédiatement
    \item Vous pouvez le modifier si nécessaire
\end{enumerate}

\astuce{Les thèmes exportés peuvent être partagés avec d'autres utilisateurs de l'application. Idéal pour standardiser l'apparence des attestations dans plusieurs services.}

\subsection{Support bilingue français/anglais}

\subsubsection{Activer le mode anglais}

\begin{enumerate}[leftmargin=*]
    \item Assurez-vous que votre fichier Excel contient les colonnes de traduction (\texttt{*\_EN})
    \item Dans l'onglet \textbf{Paramètres}, sélectionnez \champ{Langue : Anglais}
    \item Générez normalement
    \item Les attestations seront en anglais avec :
    \begin{itemize}
        \item Titre : "Certificate of Achievement" au lieu de "Attestation de Réussite"
        \item Contenu traduit automatiquement
        \item Valeurs des colonnes \texttt{*\_EN} utilisées
    \end{itemize}
\end{enumerate}

\subsubsection{Exemple de traduction}

\begin{table}[h]
\centering
\begin{tabular}{|l|l|}
\hline
\textbf{Français} & \textbf{Anglais (avec colonnes \_EN)} \\
\hline
Domaine: Sciences de la Santé & Field: Health Sciences \\
\hline
Parcours: Médecine Générale & Program: General Medicine \\
\hline
Spécialité: Cardiologie & Specialty: Cardiology \\
\hline
Mention: Très Bien & Grade: Excellent \\
\hline
\end{tabular}
\caption{Exemple de traduction automatique}
\end{table}

\newpage
